\chapter{Double Integrals}

A double integral $\iint_R f(x,y)\,\mathrm{d}A$ measures the signed
volume between the surface $z=f(x,y)$ and the region $R\subseteq\R^{2}$.
The same object computes total quantities over a region, total
production from a spatial intensity, the integral of a joint probability
density, an expected value of a function of two random variables.

\section{Core Rules}

\begin{rulebox}{Iterated integrals over a rectangle}
If $R=[a,b]\times[c,d]$ and $f$ is continuous, then by Fubini's theorem
\[
  \iint_{R} f(x,y)\,\mathrm{d}A
    \;=\; \int_{a}^{b}\int_{c}^{d} f(x,y)\,\mathrm{d}y\,\mathrm{d}x
    \;=\; \int_{c}^{d}\int_{a}^{b} f(x,y)\,\mathrm{d}x\,\mathrm{d}y .
\]
The inner integral is computed first, treating the other variable as a
constant.
\end{rulebox}

\begin{rulebox}{Separable integrand on a rectangle}
If $f(x,y)=g(x)h(y)$, then
\[
  \iint_{R} g(x)h(y)\,\mathrm{d}A
    \;=\; \left(\int_{a}^{b} g(x)\,\mathrm{d}x\right)\left(\int_{c}^{d} h(y)\,\mathrm{d}y\right) .
\]
\end{rulebox}

\begin{rulebox}{Type I region}
If $R=\{(x,y): a\le x\le b,\; g_{1}(x)\le y\le g_{2}(x)\}$, then
\[
  \iint_{R} f(x,y)\,\mathrm{d}A
    \;=\; \int_{a}^{b}\int_{g_{1}(x)}^{g_{2}(x)} f(x,y)\,\mathrm{d}y\,\mathrm{d}x .
\]
\end{rulebox}

\begin{rulebox}{Type II region}
If $R=\{(x,y): c\le y\le d,\; h_{1}(y)\le x\le h_{2}(y)\}$, then
\[
  \iint_{R} f(x,y)\,\mathrm{d}A
    \;=\; \int_{c}^{d}\int_{h_{1}(y)}^{h_{2}(y)} f(x,y)\,\mathrm{d}x\,\mathrm{d}y .
\]
\end{rulebox}

\begin{rulebox}{Changing order of integration}
If the bounds for one order are awkward (or the inner antiderivative is
not elementary), redraw the region and re-express it in the other type.
The integral itself is unchanged.
\end{rulebox}

\begin{rulebox}{Joint densities}
If $(X,Y)$ has joint pdf $f(x,y)\ge 0$ on $\R^{2}$ with
$\iint_{\R^{2}} f=1$, then
\[
  P\bigl((X,Y)\in A\bigr)=\iint_{A} f(x,y)\,\mathrm{d}A,\qquad
  \mathbb{E}[g(X,Y)]=\iint_{\R^{2}} g(x,y)f(x,y)\,\mathrm{d}A .
\]
\end{rulebox}

\section{Practice Problems}

\begin{problems}
  \item $\displaystyle\int_{0}^{1}\int_{0}^{2} (x+y)\,\mathrm{d}y\,\mathrm{d}x$
  \item $\displaystyle\int_{0}^{1}\int_{0}^{1} xy\,\mathrm{d}y\,\mathrm{d}x$
  \item $\displaystyle\int_{0}^{2}\int_{1}^{3} x^{2}y\,\mathrm{d}y\,\mathrm{d}x$
  \item $\displaystyle\int_{0}^{1}\int_{0}^{1} \mathrm{e}^{x+y}\,\mathrm{d}y\,\mathrm{d}x$
  \item $\displaystyle\int_{0}^{1}\int_{0}^{1} (x^{2}+y^{2})\,\mathrm{d}y\,\mathrm{d}x$
  \item Switch the order:
        $\displaystyle\int_{0}^{1}\int_{0}^{2} (x+y)\,\mathrm{d}y\,\mathrm{d}x$
        (verify Fubini)
  \item $\displaystyle\iint_{R} 1\,\mathrm{d}A$ where $R$ is the triangle
        with vertices $(0,0)$, $(1,0)$, $(0,1)$.
  \item $\displaystyle\iint_{R} x\,\mathrm{d}A$ over the same triangle.
  \item $\displaystyle\iint_{R} xy\,\mathrm{d}A$ over the same triangle.
  \item $\displaystyle\int_{0}^{1}\int_{0}^{x} y\,\mathrm{d}y\,\mathrm{d}x$
        \quad (Type I)
  \item $\displaystyle\int_{0}^{1}\int_{x}^{1} y\,\mathrm{d}y\,\mathrm{d}x$,
        then switch the order and recompute.
  \item $\displaystyle\int_{0}^{1}\int_{y^{2}}^{1} 1\,\mathrm{d}x\,\mathrm{d}y$
        \quad (Type II)
  \item $\displaystyle\int_{0}^{1}\int_{0}^{x} \mathrm{e}^{x}\,\mathrm{d}y\,\mathrm{d}x$
  \item Region $R=[0,1]\times[0,1]$. Compute
        $\displaystyle\iint_{R} \frac{1}{(1+x)(1+y)}\,\mathrm{d}A$.
  \item Region $R=[0,\pi]\times[0,1]$. Compute
        $\displaystyle\iint_{R} y\cos x\,\mathrm{d}A$.
  \item Switch the order of integration:
        $\displaystyle\int_{0}^{1}\int_{x}^{1} \mathrm{e}^{y^{2}}\,\mathrm{d}y\,\mathrm{d}x$,
        then evaluate.
  \item Joint pdf $f(x,y)=4xy$ on $[0,1]^{2}$ (and $0$ elsewhere).
        Verify $\iint f=1$.
  \item For the pdf above, compute $\mathbb{E}[XY]$.
  \item Volume under $z=6-x-y$ over $R=[0,1]\times[0,1]$.
  \item Average value of $f(x,y)=x+y$ over $[0,1]\times[0,1]$.
  \item Joint pdf $f(x,y)=\mathrm{e}^{-(x+y)}$ on $(0,\infty)^{2}$.
        Verify $\iint f=1$.
  \item For that pdf, compute $\mathbb{E}[X+Y]$.
  \item Compute $\displaystyle\iint_{R} xy\,\mathrm{d}A$ where $R$ is
        bounded by $y=x$, $y=x^{2}$, $x\in[0,1]$.
  \item Region $R$ bounded by $y=0$, $y=1$, $x=y^{2}$, $x=1$. Compute
        $\displaystyle\iint_{R} x\,\mathrm{d}A$.
  \item Iterated production: two inputs $L\in[0,4]$, $K\in[0,4]$.
        Average output of $Q(L,K)=L^{1/2}K^{1/2}$ over the rectangle.
\end{problems}

\section{Solutions}

\begin{solutions}
  \item Inner: $\int_{0}^{2}(x+y)\,\mathrm{d}y=2x+2$. Outer:
        $\int_{0}^{1}(2x+2)\,\mathrm{d}x=1+2=3$.

  \item Separable:
        $\bigl(\int_{0}^{1} x\,\mathrm{d}x\bigr)\bigl(\int_{0}^{1} y\,\mathrm{d}y\bigr)
        =\frac{1}{2}\cdot\frac{1}{2}=\frac{1}{4}$.

  \item Separable:
        $\bigl(\int_{0}^{2} x^{2}\,\mathrm{d}x\bigr)\bigl(\int_{1}^{3} y\,\mathrm{d}y\bigr)
        =\frac{8}{3}\cdot 4=\frac{32}{3}$.

  \item Separable: $\mathrm{e}^{x+y}=\mathrm{e}^{x}\mathrm{e}^{y}$.
        $\bigl(\int_{0}^{1} \mathrm{e}^{x}\,\mathrm{d}x\bigr)^{2}=(\mathrm{e}-1)^{2}$.

  \item $\int_{0}^{1}\bigl(x^{2}\cdot 1+\tfrac13\bigr)\,\mathrm{d}x
        =\tfrac13+\tfrac13=\tfrac23$. (Inner antiderivative w.r.t.\ $y$:
        $x^{2}y+y^{3}/3$, evaluated at $y=1$.)

  \item Other order:
        $\int_{0}^{2}\int_{0}^{1}(x+y)\,\mathrm{d}x\,\mathrm{d}y
        =\int_{0}^{2}\bigl(\tfrac12+y\bigr)\,\mathrm{d}y=1+2=3$. Same
        answer, confirming Fubini.

  \item Triangle has area $1/2$. Via integration:
        $\int_{0}^{1}\int_{0}^{1-x} 1\,\mathrm{d}y\,\mathrm{d}x
        =\int_{0}^{1}(1-x)\,\mathrm{d}x=\frac{1}{2}$.

  \item $\int_{0}^{1}\int_{0}^{1-x} x\,\mathrm{d}y\,\mathrm{d}x
        =\int_{0}^{1} x(1-x)\,\mathrm{d}x=\frac{1}{2}-\frac{1}{3}=\frac{1}{6}$.

  \item $\int_{0}^{1}\int_{0}^{1-x} xy\,\mathrm{d}y\,\mathrm{d}x
        =\int_{0}^{1} x\cdot\frac{(1-x)^{2}}{2}\,\mathrm{d}x
        =\frac{1}{2}\int_{0}^{1} (x-2x^{2}+x^{3})\,\mathrm{d}x
        =\frac{1}{2}\bigl(\tfrac12-\tfrac23+\tfrac14\bigr)
        =\frac{1}{2}\cdot\frac{1}{12}=\frac{1}{24}$.

  \item Inner: $\int_{0}^{x} y\,\mathrm{d}y=\frac{x^{2}}{2}$. Outer:
        $\int_{0}^{1} \frac{x^{2}}{2}\,\mathrm{d}x=\frac{1}{6}$.

  \item Inner: $\int_{x}^{1} y\,\mathrm{d}y=\frac{1-x^{2}}{2}$. Outer:
        $\int_{0}^{1} \frac{1-x^{2}}{2}\,\mathrm{d}x
        =\frac{1}{2}\bigl(1-\tfrac13\bigr)=\frac{1}{3}$. Other order:
        region is $\{(x,y): 0\le y\le 1,\, 0\le x\le y\}$.
        $\int_{0}^{1}\int_{0}^{y} y\,\mathrm{d}x\,\mathrm{d}y
        =\int_{0}^{1} y^{2}\,\mathrm{d}y=\frac{1}{3}$. \checkmark

  \item Inner: $\int_{y^{2}}^{1} 1\,\mathrm{d}x=y-y^{2}$. Outer:
        $\int_{0}^{1} (y-y^{2})\,\mathrm{d}y=\frac{1}{2}-\frac{1}{3}=\frac{1}{6}$.

  \item Inner: $\int_{0}^{x} \mathrm{e}^{x}\,\mathrm{d}y=x\mathrm{e}^{x}$
        (constant w.r.t.\ $y$). Outer:
        $\int_{0}^{1} x\mathrm{e}^{x}\,\mathrm{d}x
        =(x-1)\mathrm{e}^{x}\big|_{0}^{1}=0-(-1)=1$.

  \item Separable. $\int_{0}^{1} \frac{1}{1+x}\,\mathrm{d}x=\ln 2$. By
        symmetry the $y$-integral is $\ln 2$ too. Product: $(\ln 2)^{2}$.

  \item Separable. $\int_{0}^{\pi} \cos x\,\mathrm{d}x=0$. So the whole
        integral is $0$ (regardless of the $y$-piece).

  \item Region: $0\le x\le y\le 1$. Switch:
        $\int_{0}^{1}\int_{0}^{y} \mathrm{e}^{y^{2}}\,\mathrm{d}x\,\mathrm{d}y
        =\int_{0}^{1} y\,\mathrm{e}^{y^{2}}\,\mathrm{d}y
        =\frac{1}{2}(\mathrm{e}-1)$.

  \item $\int_{0}^{1}\int_{0}^{1} 4xy\,\mathrm{d}y\,\mathrm{d}x
        =4\cdot\frac{1}{2}\cdot\frac{1}{2}=1$. \checkmark

  \item $\mathbb{E}[XY]=\int_{0}^{1}\int_{0}^{1} xy\cdot 4xy\,\mathrm{d}y\,\mathrm{d}x
        =4\bigl(\int_{0}^{1} x^{2}\,\mathrm{d}x\bigr)\bigl(\int_{0}^{1} y^{2}\,\mathrm{d}y\bigr)
        =4\cdot\frac{1}{3}\cdot\frac{1}{3}=\frac{4}{9}$.

  \item $\int_{0}^{1}\int_{0}^{1} (6-x-y)\,\mathrm{d}y\,\mathrm{d}x
        =\int_{0}^{1} \bigl(6-x-\tfrac12\bigr)\,\mathrm{d}x
        =\bigl(\tfrac{11}{2}-\tfrac12\bigr)=5$.

  \item Area of the rectangle is $1$; the integral
        $\iint (x+y)\,\mathrm{d}A=\frac{1}{2}+\frac{1}{2}=1$. Average
        $=1$.

  \item $\int_{0}^{\infty} \mathrm{e}^{-x}\,\mathrm{d}x=1$, and the same
        in $y$. By separability, $\iint f=1$. \checkmark

  \item $X,Y$ are i.i.d.\ $\mathrm{Exp}(1)$, so $\mathbb{E}[X]=\mathbb{E}[Y]=1$
        and $\mathbb{E}[X+Y]=2$. Check directly:
        $\int_{0}^{\infty}\int_{0}^{\infty} (x+y)\mathrm{e}^{-x}\mathrm{e}^{-y}\,\mathrm{d}y\,\mathrm{d}x
        =\int_{0}^{\infty} \bigl(x\mathrm{e}^{-x}\cdot 1+\mathrm{e}^{-x}\cdot 1\bigr)\,\mathrm{d}x
        =1+1=2$. \checkmark

  \item On $[0,1]$, $x\ge x^{2}$, so
        $R=\{(x,y): 0\le x\le 1,\, x^{2}\le y\le x\}$.
        \[
          \int_{0}^{1}\int_{x^{2}}^{x} xy\,\mathrm{d}y\,\mathrm{d}x
            =\int_{0}^{1} x\cdot\frac{y^{2}}{2}\Big|_{x^{2}}^{x}\,\mathrm{d}x
            =\frac{1}{2}\int_{0}^{1} x(x^{2}-x^{4})\,\mathrm{d}x
            =\frac{1}{2}\int_{0}^{1} (x^{3}-x^{5})\,\mathrm{d}x
            =\frac{1}{2}\bigl(\tfrac14-\tfrac16\bigr)=\frac{1}{24} .
        \]

  \item Type II: $0\le y\le 1$, $y^{2}\le x\le 1$.
        \[
          \int_{0}^{1}\int_{y^{2}}^{1} x\,\mathrm{d}x\,\mathrm{d}y
            =\int_{0}^{1} \frac{1-y^{4}}{2}\,\mathrm{d}y
            =\frac{1}{2}\bigl(1-\tfrac15\bigr)=\frac{2}{5} .
        \]

  \item Separable. $\int_{0}^{4} L^{1/2}\,\mathrm{d}L
        =\bigl[\tfrac23 L^{3/2}\bigr]_{0}^{4}=\frac{16}{3}$. Same for
        $K$. So $\iint Q\,\mathrm{d}A=\frac{16}{3}\cdot\frac{16}{3}=\frac{256}{9}$.
        Area of the rectangle is $16$, so the average is
        $\dfrac{256}{9\cdot 16}=\dfrac{16}{9}$.
\end{solutions}

\section{Additional Problems}

\begin{problems}
  \item $\displaystyle\int_{0}^{2}\int_{0}^{1} (2x+y)\,\mathrm{d}y\,\mathrm{d}x$
  \item $\displaystyle\int_{0}^{1}\int_{0}^{2} x^{2}y\,\mathrm{d}y\,\mathrm{d}x$
  \item $\displaystyle\int_{0}^{1}\int_{0}^{1} (x+y)^{2}\,\mathrm{d}y\,\mathrm{d}x$
  \item $\displaystyle\int_{0}^{1}\int_{1}^{2} \mathrm{e}^{x-y}\,\mathrm{d}y\,\mathrm{d}x$
  \item $\displaystyle\int_{0}^{2}\int_{0}^{3} (x^{2}+y^{2})\,\mathrm{d}y\,\mathrm{d}x$
  \item $\displaystyle\iint_{R} xy\,\mathrm{d}A$ where $R$ is the
        triangle $(0,0)$, $(2,0)$, $(0,2)$.
  \item $\displaystyle\iint_{R} y\,\mathrm{d}A$ where $R$ is the
        triangle $(0,0)$, $(1,0)$, $(1,1)$.
  \item $\displaystyle\int_{0}^{2}\int_{x^{2}}^{2x} 1\,\mathrm{d}y\,\mathrm{d}x$
  \item $\displaystyle\int_{0}^{1}\int_{0}^{\sqrt{y}} 2y\,\mathrm{d}x\,\mathrm{d}y$
        \quad (Type II)
  \item $\displaystyle\iint_{R} (x+y)\,\mathrm{d}A$ where $R$ is
        bounded by $y=x$, $y=2x$, $x\in[0,1]$.
  \item $\displaystyle\int_{0}^{1}\int_{0}^{2} x\,\mathrm{e}^{y}\,\mathrm{d}y\,\mathrm{d}x$
  \item $\displaystyle\iint_{R} \frac{1}{(x+y+1)^{2}}\,\mathrm{d}A$
        over $R=[0,1]^{2}$.
  \item Switch the order of integration:
        $\displaystyle\int_{0}^{1}\int_{\sqrt{x}}^{1} \mathrm{e}^{y^{3}}\,\mathrm{d}y\,\mathrm{d}x$,
        then evaluate.
  \item Joint pdf $f(x,y)=6\,\mathrm{e}^{-2x-3y}$ on $(0,\infty)^{2}$.
        Verify $\iint f=1$ and compute $\mathbb{E}[X+Y]$.
  \item Cobb--Douglas production $Q(L,K)=L^{1/3}K^{2/3}$ on
        $[0,8]\times[0,8]$. Compute the average output.
\end{problems}

\section{Additional Solutions}

\begin{solutions}
  \item Inner: $\int_{0}^{1}(2x+y)\,\mathrm{d}y=2x+\tfrac12$. Outer:
        $\int_{0}^{2}\bigl(2x+\tfrac12\bigr)\,\mathrm{d}x=4+1=5$.

  \item Separable:
        $\bigl(\int_{0}^{1} x^{2}\,\mathrm{d}x\bigr)\bigl(\int_{0}^{2} y\,\mathrm{d}y\bigr)
        =\frac{1}{3}\cdot 2=\frac{2}{3}$.

  \item Inner: $\int_{0}^{1}(x^{2}+2xy+y^{2})\,\mathrm{d}y
        =x^{2}+x+\tfrac13$. Outer:
        $\int_{0}^{1}\bigl(x^{2}+x+\tfrac13\bigr)\,\mathrm{d}x
        =\tfrac13+\tfrac12+\tfrac13=\frac{7}{6}$.

  \item Separable: $\mathrm{e}^{x-y}=\mathrm{e}^{x}\cdot\mathrm{e}^{-y}$.
        $\bigl(\int_{0}^{1} \mathrm{e}^{x}\,\mathrm{d}x\bigr)
        \bigl(\int_{1}^{2} \mathrm{e}^{-y}\,\mathrm{d}y\bigr)
        =(\mathrm{e}-1)(\mathrm{e}^{-1}-\mathrm{e}^{-2})
        =(\mathrm{e}-1)^{2}/\mathrm{e}^{2}$.

  \item Inner: $\int_{0}^{3}(x^{2}+y^{2})\,\mathrm{d}y=3x^{2}+9$. Outer:
        $\int_{0}^{2}(3x^{2}+9)\,\mathrm{d}x=8+18=26$.

  \item Type I: $0\le x\le 2$, $0\le y\le 2-x$.
        \[
          \int_{0}^{2}\int_{0}^{2-x} xy\,\mathrm{d}y\,\mathrm{d}x
            =\int_{0}^{2} x\cdot\frac{(2-x)^{2}}{2}\,\mathrm{d}x
            =\frac{1}{2}\int_{0}^{2} (4x-4x^{2}+x^{3})\,\mathrm{d}x
            =\frac{1}{2}\bigl(8-\tfrac{32}{3}+4\bigr)=\frac{2}{3} .
        \]

  \item Type I: $0\le x\le 1$, $0\le y\le x$.
        \[
          \int_{0}^{1}\int_{0}^{x} y\,\mathrm{d}y\,\mathrm{d}x
            =\int_{0}^{1} \frac{x^{2}}{2}\,\mathrm{d}x=\frac{1}{6} .
        \]

  \item Curves $y=x^{2}$ and $y=2x$ meet at $x=0$ and $x=2$, with
        $2x\ge x^{2}$ on $[0,2]$. Inner: $2x-x^{2}$. Outer:
        \[
          \int_{0}^{2} (2x-x^{2})\,\mathrm{d}x=4-\frac{8}{3}=\frac{4}{3} .
        \]

  \item Inner: $\int_{0}^{\sqrt{y}} 2y\,\mathrm{d}x=2y\sqrt{y}=2y^{3/2}$.
        Outer: $\int_{0}^{1} 2y^{3/2}\,\mathrm{d}y=2\cdot\frac{2}{5}=\frac{4}{5}$.

  \item Type I: $0\le x\le 1$, $x\le y\le 2x$. Inner:
        $\int_{x}^{2x}(x+y)\,\mathrm{d}y
        =\bigl[xy+\tfrac{y^{2}}{2}\bigr]_{x}^{2x}
        =(2x^{2}+2x^{2})-\bigl(x^{2}+\tfrac{x^{2}}{2}\bigr)
        =4x^{2}-\tfrac{3x^{2}}{2}=\tfrac{5x^{2}}{2}$.
        \[
          \int_{0}^{1} \frac{5x^{2}}{2}\,\mathrm{d}x=\frac{5}{6} .
        \]

  \item Separable:
        $\bigl(\int_{0}^{1} x\,\mathrm{d}x\bigr)\bigl(\int_{0}^{2} \mathrm{e}^{y}\,\mathrm{d}y\bigr)
        =\frac{1}{2}(\mathrm{e}^{2}-1)$.

  \item Inner ($y$ first): $\int_{0}^{1}(x+y+1)^{-2}\,\mathrm{d}y
        =\bigl[-\tfrac{1}{x+y+1}\bigr]_{0}^{1}
        =-\tfrac{1}{x+2}+\tfrac{1}{x+1}=\dfrac{1}{(x+1)(x+2)}$. Outer
        (partial fractions $1/((x+1)(x+2))=1/(x+1)-1/(x+2)$):
        \[
          \int_{0}^{1}\Bigl(\frac{1}{x+1}-\frac{1}{x+2}\Bigr)\,\mathrm{d}x
            =\bigl[\ln(x+1)-\ln(x+2)\bigr]_{0}^{1}
            =(\ln 2-\ln 3)-(0-\ln 2)
            =2\ln 2-\ln 3=\ln\frac{4}{3} .
        \]

  \item Region: $0\le x\le 1$, $\sqrt{x}\le y\le 1$, equivalently
        $0\le y\le 1$, $0\le x\le y^{2}$. After switching,
        \[
          \int_{0}^{1}\int_{0}^{y^{2}} \mathrm{e}^{y^{3}}\,\mathrm{d}x\,\mathrm{d}y
            =\int_{0}^{1} y^{2}\mathrm{e}^{y^{3}}\,\mathrm{d}y
            =\frac{1}{3}(\mathrm{e}-1) ,
        \]
        using $u=y^{3}$.

  \item Separable:
        $\int_{0}^{\infty} 6\,\mathrm{e}^{-2x-3y}\,\mathrm{d}x\,\mathrm{d}y
        =6\cdot\frac{1}{2}\cdot\frac{1}{3}=1$. \checkmark\
        Marginals: $X\sim\mathrm{Exp}(2)$, $Y\sim\mathrm{Exp}(3)$. So
        $\mathbb{E}[X]=1/2$, $\mathbb{E}[Y]=1/3$, and
        $\mathbb{E}[X+Y]=5/6$.

  \item Separable. $\int_{0}^{8} L^{1/3}\,\mathrm{d}L
        =\bigl[\tfrac34 L^{4/3}\bigr]_{0}^{8}=\frac{3}{4}\cdot 16=12$.
        $\int_{0}^{8} K^{2/3}\,\mathrm{d}K
        =\bigl[\tfrac35 K^{5/3}\bigr]_{0}^{8}=\frac{3}{5}\cdot 32=\frac{96}{5}$.
        Total $\iint Q\,\mathrm{d}A=12\cdot\frac{96}{5}=\frac{1152}{5}$.
        Area $=64$, so average
        $=\dfrac{1152}{5\cdot 64}=\dfrac{18}{5}=3.6$.
\end{solutions}
